\chapter{Tools}
\label{ch:tools}

Tools are typed functions with LLM-visible descriptions. The compiler auto-generates JSON schemas from the type signature, enabling agents to call tools with structured arguments.

\section{Tool Declaration}

\begin{lstlisting}
tool tool_name(param: Type, param: Type = default) -> ReturnType {
    """Description visible to the LLM"""

    // implementation
}
\end{lstlisting}

A tool is like \keyword{fn} with two additions:
\begin{enumerate}
    \item The \keyword{tool} keyword instead of \keyword{fn}
    \item A mandatory triple-quoted description as the first element of the body
\end{enumerate}

\section{Basic Examples}

\subsection{HTTP Tool}

\begin{lstlisting}
import "tools/http"
import "json"

struct SearchResult {
    url: string
    title: string
    snippet: string
}

tool search(query: string, max_results: int = 5) -> [SearchResult] {
    """
    Search the web for information.
    Returns up to max_results results.
    Supports boolean operators in query.
    """

    resp, err = http.get("https://api.search.com?q={query}&n={max_results}")
    if err != nil { return [], err }
    return json.decode(resp.body, [SearchResult])
}
\end{lstlisting}

\subsection{Database Tool}

\begin{lstlisting}
import "tools/postgres"

tool get_user(user_id: string) -> User {
    """Look up a user by their ID"""

    user, err = pg.query_one("SELECT * FROM users WHERE id = $1", [user_id])
    if err != nil { return User{}, err }
    return user
}
\end{lstlisting}

\subsection{External Service Tool}

\begin{lstlisting}
import "tools/http"

tool send_slack(channel: string, message: string) -> bool {
    """Send a message to a Slack channel"""

    resp, err = http.post("https://slack.com/api/chat.postMessage", {
        channel: channel,
        text: message
    })
    return err == nil
}
\end{lstlisting}

\subsection{Agent-Delegating Tool}

Tools can call agents internally:

\begin{lstlisting}
tool summarize(text: string) -> string {
    """Summarize text into key points"""

    result, _ = Writer.ask("Summarize concisely: {text}")
    return result
}
\end{lstlisting}

\section{Triple-Quote Description}

The description is mandatory and must be the first element of the tool body. It is a triple-quoted string (\code{"""..."""}) that can span multiple lines:

\begin{lstlisting}
tool analyze(data: [DataPoint]) -> Analysis {
    """
    Analyze a dataset and produce a statistical summary.

    Includes:
    - Mean, median, and standard deviation
    - Outlier detection
    - Trend analysis over time

    Use this tool when the user asks about patterns or statistics.
    """

    // implementation
}
\end{lstlisting}

\begin{note}
The compiler will emit an error if a tool is missing its description. This is enforced because the description is sent to the LLM as part of the tool-calling protocol.
\end{note}

\section{Auto-Generated JSON Schema}

The compiler generates a JSON schema from the tool's type signature. For example:

\begin{lstlisting}
tool search(query: string, max_results: int = 5) -> [SearchResult] {
    """Search the web for information"""
    // ...
}
\end{lstlisting}

Generates (conceptually):

\begin{lstlisting}[language={}]
{
  "name": "search",
  "description": "Search the web for information",
  "parameters": {
    "type": "object",
    "properties": {
      "query": { "type": "string" },
      "max_results": { "type": "integer", "default": 5 }
    },
    "required": ["query"]
  }
}
\end{lstlisting}

This happens at compile time. No hand-written JSON schemas are needed.

\section{Tool Packs (Standard Library)}

\haira{} ships with built-in tool packs for common integrations:

\begin{lstlisting}
import "tools/http"       // http.get, http.post, http.put, http.delete
import "tools/slack"      // slack.send, slack.read_channel
import "tools/github"     // github.create_issue, github.list_prs, github.get_diff
import "tools/postgres"   // pg.query, pg.query_one, pg.insert, pg.update
import "tools/email"      // email.send
import "tools/fs"         // fs.read, fs.write, fs.list
\end{lstlisting}

Tool packs provide pre-built \keyword{tool} declarations. They can be used directly in agent \code{tools} lists or called from within other tools and workflows.

\section{Error Handling in Tools}

Tools follow \haira{}'s standard error handling pattern:

\begin{lstlisting}
tool fetch_page(url: string) -> string {
    """Fetch the content of a web page"""

    resp, err = http.get(url)
    if err != nil {
        return "", err
    }
    if resp.status != 200 {
        return "", Error{message: "HTTP {resp.status}"}
    }
    return resp.body
}
\end{lstlisting}

When a tool returns an error, the agent receives the error message and can decide how to proceed (retry, try a different tool, or report the error to the user).

\section{Tool Grammar}

\begin{lstlisting}[language=EBNF]
tool_decl     = "tool" identifier "(" [ params ] ")" [ "->" type ]
                "{" triple_string { statement } "}" ;

triple_string = '"""' { any_char } '"""' ;
\end{lstlisting}
